\subsection{Related Work}

A limited number of previous works have learned dictionaries of sparsely-activating features in pre-trained models, including \citet{yun2021transformer} and \citet{sharkey2023technical}, the latter of which motivated this work. However, similar methods have been applied in other domains, in particular in understanding neurons in the visual cortex \citep{olshausen2004sparse, wright2022high}. 

In contrast to our approach, where we try to impose sparsity after training, many previous works have encouraged sparsity in neural networks via changes to the architecture or training process. These approaches include altering the attention mechanism \citep{correia2019adaptively}, adding $\ell^1$ penalties to neuron activations \citep{kasioumis2021elite, georgiadis2019accelerating}, pruning neurons \citep{frankle2018lottery}, and using the softmax function as the non-linearity in the MLP layers \citep{elhage2022solu}. However, training a state-of-the-art foundation model with these additional constraints is difficult \citep{elhage2022solu}, and improvements to interpretability are not always realized \citep{Meister2021IsSA}.

\subsection{Conclusion}

Sparse autoencoders are a scalable, unsupervised approach to disentangling language model network features from superposition. Our approach requires only unlabelled model activations and uses orders of magnitude less compute than the training of the original models. We have demonstrated that the dictionary features we learn are more interpretable by autointerpretation and are more monosemantic than comparable methods. This approach could facilitate the mapping of model circuits, targeted model editing, and a better understanding of model representations.
